\section{国内外现状}

近年来，
随着虚拟现实、数字孪生以及游戏制作等应用的发展，
三维内容生成技术逐渐成为计算机图形学与计算机视觉领域的重要研究方向。
其中，
三维材质生成与材质迁移是三维内容自动化生产的重要组成部分，
能够为三维模型提供真实且多样的表面外观，
因此受到了广泛关注。

首先，
基于深度学习的材质生成方法成为当前的重要研究方向。
研究者通常利用卷积神经网络或生成模型，
从大量真实材质数据中学习纹理分布，
从而生成高质量的材质贴图。
随着生成式模型的发展，
扩散模型等方法被逐渐应用于三维内容生成任务，
通过引入二维生成模型的先验信息，
实现由文本或图像驱动的三维形状与材质生成，
显著提升了生成材质的细节表现能力与多样性。

其次，
神经渲染表示方法为三维材质建模与迁移提供了新的思路。
以神经辐射场（Neural Radiance Field，NeRF）为代表的方法，
通过神经网络对三维空间中的颜色和密度进行连续建模，
能够在多视角条件下学习物体的几何结构与外观信息。
在此基础上，
部分研究进一步将外观表示分解为几何与材质成分，
从而实现材质编辑与迁移，
使得模型能够在保持几何结构不变的情况下替换或调整表面材质。

此外，
基于显式三维表示的材质迁移方法也是当前研究的热点之一。
该类方法通常在网格或点云等显式结构上进行纹理映射与优化，
通过风格迁移、纹理重建或跨模型纹理映射等技术，
将一种物体的材质外观迁移到另一三维模型上，
从而实现快速的材质重用与外观编辑。
这类方法生成结果可以直接应用于传统图形渲染管线，
在游戏制作与三维内容生产中具有较高的实际应用价值。

总体而言，
当前三维材质相关研究主要集中在材质生成、材质迁移以及可编辑三维外观建模等方向。
随着生成式模型与神经渲染技术的发展，
三维材质生成与迁移的自动化程度不断提高。
然而，
在多视角一致性、真实物理材质建模以及高效可控编辑等方面，
仍然存在进一步研究的空间。